\documentclass[preview]{standalone}

\usepackage{amsmath}
\usepackage{amssymb}
\usepackage{stellar}
\usepackage{definitions}

\begin{document}

\id{psicologia-problem-solving-ragionamento}
\genpage

\section{Problem solving e ragionamento}

\begin{snippet}{problem-solving-ragionamento-introduzione}
    Problem solving e ragionamento combinano informazioni contingenti con
    conoscenze immagazzinate in memoria. Lo scopo è raggiungere una conclusione
    o trovare una soluzione quando la situazione attuale non coincide con quella
    desiderata.
\end{snippet}

\begin{snippetdefinition}{problem-solving-definition}{Problem solving}
    Il \emph{problem solving} è l'attività cognitiva attraverso cui si cerca di
    ridurre la discrepanza tra ciò che si sa o si ha nello stato attuale e ciò
    che si vuole sapere o ottenere nello stato desiderato.
\end{snippetdefinition}

\begin{snippetexample}{nove-punti-example}{Problema dei nove punti}
    Nel problema dei nove punti bisogna unire tutti i punti con quattro linee
    rette connesse senza staccare la penna dal foglio. La difficoltà deriva
    dalla tendenza a rappresentare implicitamente i punti come racchiusi entro
    un quadrato, mentre la soluzione richiede di uscire da tale cornice mentale.
\end{snippetexample}

\includesnpt[width=45\%,src=/snippet/static-psychology/nine-dot-problem.jpg]{centered-img}
\includesnpt[width=45\%,src=/snippet/static-psychology/nine-dot-solution.jpg]{centered-img}

\begin{snippetdefinition}{spazio-problema-definition}{Spazio del problema}
    Lo \emph{spazio del problema} è l'insieme formato da stato iniziale, stato
    finale e operatori disponibili per passare dall'uno all'altro.
\end{snippetdefinition}

\begin{snippetdefinition}{problema-ben-definito-definition}{Problema ben definito}
    Un \emph{problema ben definito} è un problema in cui stato iniziale, stato
    finale e operatori sono chiaramente specificati. Il compito consiste nello
    scoprire come usare le operazioni disponibili per raggiungere la soluzione.
\end{snippetdefinition}

\begin{snippetdefinition}{algoritmo-definition}{Algoritmo}
    Un \emph{algoritmo} è una procedura passo dopo passo che fornisce sempre la
    soluzione corretta a un particolare tipo di problema.
\end{snippetdefinition}

\begin{snippetdefinition}{euristica-problem-solving-definition}{Euristica}
    Un'\emph{euristica} è una strategia cognitiva o regola empirica usata come
    scorciatoia nella soluzione di problemi complessi.
\end{snippetdefinition}

\begin{snippetexample}{euristica-romanzo-giallo-example}{Euristica nel romanzo giallo}
    Leggendo un romanzo giallo, si potrebbe escludere il maggiordomo come
    assassino usando l'euristica secondo cui l'autore non sceglierebbe mai una
    trama così scontata.
\end{snippetexample}

\begin{snippetdefinition}{fissita-funzionale-definition}{Fissità funzionale}
    La \emph{fissità funzionale} è un blocco mentale che ostacola la risoluzione
    dei problemi, impedendo di vedere una nuova funzione per un oggetto
    precedentemente associato a un altro scopo.
\end{snippetdefinition}

\begin{snippetexample}{problema-candela-example}{Problema della candela}
    Nel problema della candela si hanno una candela, una scatola di puntine e
    fiammiferi. Il compito è attaccare la candela al muro sopra un tavolo in
    modo che la cera non goccioli sul tavolo. La soluzione richiede di
    reinterpretare la scatola di puntine non come contenitore, ma come
    piattaforma da fissare al muro.
\end{snippetexample}

\includesnpt[width=50\%,src=/snippet/static-psychology/candle-problem-functional-fixedness.jpg]{centered-img}

\begin{snippetdefinition}{ragionamento-deduttivo-definition}{Ragionamento deduttivo}
    Il \emph{ragionamento deduttivo} consiste nel trarre conclusioni a partire
    da premesse, basandosi su regole logiche che procedono dal generale al
    particolare.
\end{snippetdefinition}

\begin{snippetexample}{sillogismo-visa-example}{Sillogismo della carta Visa}
    Se un ristorante accetta tutte le principali carte di credito e la Visa è
    una delle principali carte di credito, allora si può concludere che il
    ristorante accetta la Visa. La conclusione segue dalle premesse secondo una
    struttura deduttiva.
\end{snippetexample}

\begin{snippetdefinition}{bias-credenza-definition}{Bias dovuto alla credenza}
    Il \emph{bias dovuto alla credenza} è la tendenza a giudicare valide le
    conclusioni che si ritengono credibili e non valide quelle che si giudicano
    non credibili.
\end{snippetdefinition}

\begin{snippetexample}{wason-selection-task-example}{Compito di selezione di Wason}
    Nel compito di selezione di Wason vengono presentate quattro carte:
    \texttt{A}, \texttt{D}, \texttt{4}, \texttt{7}. Per verificare la regola
    ``se una carta ha una vocale da un lato, allora ha un numero pari
    dall'altro'', bisogna girare \texttt{A} e \texttt{7}. Girare \texttt{4} non
    serve, perché la regola non dice che ogni numero pari debba avere una vocale
    sull'altro lato.
\end{snippetexample}

\includesnpt[width=65\%,src=/snippet/static-psychology/wason-selection-task.jpg]{centered-img}

\begin{snippetdefinition}{ragionamento-induttivo-definition}{Ragionamento induttivo}
    Il \emph{ragionamento induttivo} è una forma di ragionamento che usa gli
    indizi disponibili per generare conclusioni probabili ma non certe.
\end{snippetdefinition}

\begin{snippetdefinition}{soluzione-problemi-analogia-definition}{Soluzione di problemi per analogia}
    La \emph{soluzione di problemi per analogia} consiste nello stabilire
    un'analogia tra la situazione corrente e una situazione pregressa, usando una
    soluzione già sperimentata in passato.
\end{snippetdefinition}

\begin{snippetdefinition}{disposizione-mentale-definition}{Disposizione mentale}
    Una \emph{disposizione mentale} è uno stato cognitivo preesistente che può
    migliorare percezione e soluzione dei problemi, ma può anche distorcere
    l'attività mentale quando le vecchie modalità di pensiero non sono più
    adeguate.
\end{snippetdefinition}

\begin{snippetdefinition}{ragionamento-abduttivo-definition}{Ragionamento abduttivo}
    Il \emph{ragionamento abduttivo} consiste nel procedere a ritroso dagli
    effetti alle cause, nel tentativo di spiegare qualcosa che è già accaduto.
    Procede per supposizioni e non garantisce conclusioni corrette.
\end{snippetdefinition}

\begin{snippetdefinition}{fissazione-attentiva-definition}{Fissazione attentiva}
    La \emph{fissazione attentiva} è la concentrazione dell'attenzione su
    aspetti parziali e limitati di ciò che è stato detto o di ciò che è accaduto.
\end{snippetdefinition}

\includesnpt[width=58\%,src=/snippet/static-psychology/deductive-inductive-abductive-reasoning.jpg]{centered-img}

\end{document}
